% ============================================================
% MATH MACROS
% ============================================================

% --- Zahlmengen (Blackboard Bold) ---
\newcommand{\R}{\mathbb{R}}
\newcommand{\C}{\mathbb{C}}
\newcommand{\N}{\mathbb{N}}
\newcommand{\Z}{\mathbb{Z}}
\newcommand{\Q}{\mathbb{Q}}

% --- Differentialoperator (aufrechtes d) ---
\newcommand{\dd}{\,\mathrm{d}}

% --- Operatornamen ---
\DeclareMathOperator{\sgn}{sgn}
\DeclareMathOperator{\grad}{grad}
\DeclareMathOperator{\divg}{div}
\DeclareMathOperator{\rot}{rot}
\newcommand{\vect}[1]{\vec{#1}}

% --- Gepaarte Begrenzer (mathtools) ---
\DeclarePairedDelimiter{\abs}{\lvert}{\rvert}
\DeclarePairedDelimiter{\norm}{\lVert}{\rVert}

% --- Pfeile (Fix für sansmath Überlappungen) ---
% Da sansmath/Carlito bei zusammengesetzten langen Pfeilen (wie \Longrightarrow) 
% unschöne Überlappungen erzeugt, biegen wir diese auf die sauberen kurzen Pfeile um.
% Das spart auf dem Cheat Sheet zudem Platz.
\renewcommand{\implies}{\;\Rightarrow\;}
\renewcommand{\iff}{\;\Leftrightarrow\;}
\renewcommand{\Longrightarrow}{\Rightarrow}
\renewcommand{\Longleftrightarrow}{\Leftrightarrow}
\renewcommand{\longrightarrow}{\rightarrow}
\renewcommand{\longleftarrow}{\leftarrow}

% --- Summen/Limes-Mode-System ---
\newcommand{\ZSFsumMode}{auto}
\newcommand{\SetZSFsumMode}[1]{\renewcommand{\ZSFsumMode}{#1}}
\newcommand{\ZSFsumSide}[2]{\sum\nolimits_{#1}^{#2}}
\newcommand{\ZSFsumStack}[2]{\sum\limits_{#1}^{#2}}
\newcommand{\ZSFsumAuto}[2]{%
  \IfStrEqCase{\ZSFsumMode}{%
    {side}{\ZSFsumSide{#1}{#2}}%
    {stack}{\ZSFsumStack{#1}{#2}}%
    {auto}{\sum_{#1}^{#2}}%
  }[\sum_{#1}^{#2}]%
}

\newcommand{\ZSFlimMode}{auto}
\newcommand{\SetZSFlimMode}[1]{\renewcommand{\ZSFlimMode}{#1}}
\newcommand{\ZSFlimAuto}[1]{%
  \IfStrEqCase{\ZSFlimMode}{%
    {side}{\lim\nolimits_{#1}}%
    {stack}{\lim\limits_{#1}}%
    {auto}{\lim_{#1}}%
  }[\lim_{#1}]%
}

% --- Semantisches Formel-Highlighting ---
% Verknuepft visuell zusammengehoerige Teile einer Formel ueber mehrere
% Umformungsschritte hinweg. Farben zentral in 40_colors_structure.tex.
\newcommand{\ZSFmhlA}[1]{{\color{MathHLPrimary}#1}}
\newcommand{\ZSFmhlB}[1]{{\color{MathHLSecondary}#1}}
\newcommand{\ZSFmhlC}[1]{{\color{MathHLTarget}#1}}
\newcommand{\ZSFmhlD}[1]{{\color{MathHLTertiary}#1}}
